\setcounter{chaptercntr}{1}

\sectionbreak \section*{
	\gostTitleFont
	\redline
	\thechaptercntr \space
	АНАЛИЗ ПРЕДМЕТНОЙ ОБЛАСТИ И ПОСТАНОВКА ЗАДАЧИ
}

\titlespace

\subsection*{
	\gostTitleFont
	\redline
	\thechaptercntr .\thesubchaptercntr \spc
	Общие сведения о сервисах заказа такси
} \addtocounter{subchaptercntr}{1}

\subtitlespace

{\gostFont
	
	\par \redline Сервисы заказа такси представляют собой цифровые платформы, обеспечивающие взаимодействие пассажиров, водителей и операторов сервиса в рамках единого информационного пространства. Основная задача такого сервиса заключается в быстром приеме заказа, определении подходящего исполнителя, построении маршрута поездки, расчете стоимости и сопровождении заказа на всех стадиях его жизненного цикла.
	
	\par \redline Современный рынок такси существенно отличается от классической модели диспетчерской службы. Если ранее ключевая роль отводилась оператору, который вручную принимал звонки и передавал заявку водителю, то теперь значительная часть операций автоматизирована. Пользователь самостоятельно формирует заказ через мобильное приложение, указывает точку подачи автомобиля, пункт назначения, способ оплаты и дополнительные пожелания. После этого серверная часть системы выполняет обработку запроса и инициирует дальнейшие бизнес-процессы [4].
	
	\par \redline В рамках предметной области можно выделить несколько основных участников. Пассажир создает заказ, отслеживает подачу автомобиля, получает информацию о водителе и оплачивает поездку. Водитель принимает или отклоняет предложенный заказ, прибывает в точку подачи, выполняет перевозку и завершает поездку. Администратор сервиса управляет тарифами, пользователями, промокампаниями, справочниками, зонами обслуживания и механизмами контроля качества. Помимо этого, в системе могут участвовать внешние сервисы: платежные шлюзы, геокодеры, картографические платформы, сервисы уведомлений и сообщений.
	
	\par \redline Заказ такси как сущность предметной области проходит ряд состояний. В типичном сценарии последовательность выглядит так: "создан"{}, "назначается"{}, "принят водителем"{}, "водитель прибыл"{}, "поездка начата"{}, "поездка завершена"{}, затем - "оплачен"{} или "отменён"{}. Переходы между состояниями сопровождаются проверкой бизнес-правил, обновлением данных в нескольких подсистемах и отправкой уведомлений заинтересованным участникам. Поэтому серверная часть приложения должна обеспечивать не только хранение текущего состояния заказа, но и надёжную координацию и валидацию переходов между состояниями.
	
	\par \redline Существенной особенностью сервисов заказа такси является их географическая привязка. Для работы платформы необходимо постоянно обрабатывать координаты пассажиров и водителей, рассчитывать расстояния, оценивать примерное время подачи автомобиля и прогнозировать продолжительность поездки. Данные задачи требуют интеграции с картографическими системами и алгоритмами маршрутизации. При этом серверная часть должна работать с координатами в режиме, близком к реальному времени, так как даже небольшая задержка может привести к неверному выбору водителя или ухудшению пользовательского опыта.
	
	\par \redline Важной характеристикой данной предметной области является и зависимость качества сервиса от скорости обмена данными между клиентской и серверной частями. Пассажир ожидает быстрого подтверждения заказа и получения достоверной информации о подаче автомобиля. Водителю, в свою очередь, необходим своевременный доступ к сведениям о новом заказе, маршруте и статусе поездки. Следовательно, серверная часть должна обеспечивать непрерывный обмен данными и согласованное представление состояния заказа для всех участников процесса.
	
	\par \redline Еще одной характерной чертой является высокая событийность предметной области. В течение короткого промежутка времени в системе могут происходить десятки тысяч однотипных действий: создание заказов, обновление геопозиции, смена статусов поездок, расчет стоимости, отправка уведомлений, авторизация пользователей, фиксация платежных операций. Эти события возникают асинхронно, пересекаются по времени и затрагивают разные части системы. По этой причине серверная часть должна проектироваться как основа для обработки большого числа конкурентных запросов и событий.
	
	\par \redline Бизнес-логика приложений "Такси"{} не сводится только к поиску ближайшего автомобиля. На практике сервер должен поддерживать тарифные планы, динамическое ценообразование, валидацию промокодов, привязку банковских карт, историю поездок, рейтинговую систему, блокировки пользователей, обработку жалоб, возвраты средств и формирование аналитики. Дополнительно возникают задачи, связанные с безопасностью поездки, подтверждением личности водителя, хранением документов и журналированием действий пользователей. Таким образом, предметная область сочетает в себе элементы транспортной логистики, финансовых расчетов, геоинформационной обработки и управления пользовательскими аккаунтами.
	
	\par \redline С точки зрения программной реализации сервис заказа такси обычно включает несколько клиентских приложений и единую серверную платформу. К клиентским компонентам относятся мобильное приложение пассажира, мобильное приложение водителя, веб-панель администратора, а иногда и интерфейс операторов контакт-центра. Серверная часть выступает центральным звеном ИС, поскольку именно она реализует доменную логику, управляет доступом к данным, координирует взаимодействие между подсистемами и обеспечивает интеграцию с внешними платформами.
	
	\par \redline Типичный процесс выполнения заказа можно представить следующим образом. Пользователь проходит аутентификацию, отправляет запрос на создание поездки, после чего сервер валидирует входные данные и рассчитывает предварительную стоимость. Далее подсистема распределения заказов определяет список доступных водителей в заданном радиусе, выполняет ранжирование кандидатов и инициирует механизм назначения. После подтверждения заказа водителем система переводит заказ в активное состояние, сохраняет события маршрута, фиксирует время прибытия и поддерживает обмен статусами между участниками. По завершении поездки выполняется окончательный расчет стоимости, списание средств, обновление балансов и запись информации в историю пользователя.
	
	\par \redline На практике отдельные этапы жизненного цикла заказа могут различаться в зависимости от бизнес-правил конкретного сервиса. Например, может использоваться предварительное резервирование суммы поездки, автоматическое назначение водителя без подтверждения либо ручная отмена поездки со стороны оператора. Несмотря на различия в деталях реализации, базовая логика остается общей: система должна корректно принимать заказ, назначать исполнителя, сопровождать поездку и фиксировать результат выполнения услуги.
	
	\par \redline Серверная часть занимает центральное место в структуре программного комплекса. Она обеспечивает взаимодействие клиентских приложений с подсистемами авторизации, обработки заказов, геолокации, уведомлений и хранения данных. Такое положение определяет ее роль как координирующего элемента всей ИС и задает требования к производительности, сопровождаемости и возможности дальнейшего расширения.
	
	\par \redline Анализ предметной области показывает, что серверная часть приложения "Такси"{} является критически важным компонентом всей платформы. Именно от ее архитектуры зависят скорость обработки заказов, корректность бизнес-операций и удобство дальнейшего развития продукта. Следовательно, разработка серверной части должна учитывать как специфику бизнес-процессов такси, так и технические ограничения распределенных систем, даже если на текущем этапе реализуется минимально жизнеспособный продукт.
	
	\par
}

\subtitlespace

\subsection*{
	\gostTitleFont
	\redline
	\thechaptercntr .\thesubchaptercntr \spc
	Особенности разработки высоконагруженных серверных систем
} \addtocounter{subchaptercntr}{1}

\subtitlespace

{\gostFont
	
	\par \redline Высоконагруженной серверной системой принято считать программную систему, которая должна стабильно функционировать при большом количестве одновременных запросов, интенсивном обмене данными и строгих требованиях к времени отклика [2]. Для приложений заказа такси такая характеристика является естественной, поскольку нагрузка в системе распределена неравномерно и может резко возрастать в часы пик, во время неблагоприятных погодных условий, в праздничные дни и при проведении массовых мероприятий.
	
	\par \redline Одной из ключевых проблем разработки высоконагруженной системы является обеспечение горизонтальной масштабируемости. Если сервис строится таким образом, что увеличение числа пользователей требует только наращивания ресурсов одного сервера, то со временем возникает аппаратный и организационный предел. Более практичным подходом является проектирование компонентов, которые можно тиражировать на несколько экземпляров и распределять между ними входящий поток запросов. Для серверной части приложения "Такси"{} это особенно важно, так как подсистемы аутентификации, управления заказами, работы с геолокацией и уведомлениями могут испытывать различный профиль нагрузки.
	
	\par \redline Существенное значение имеет производительность операций чтения и записи. В системе такси постоянно обновляются координаты водителей, создаются новые заказы, фиксируются изменения статусов и выполняются платежные транзакции. При этом часть данных должна быть доступна почти мгновенно, например текущие координаты автомобилей или состояние активного заказа. Другие данные могут обрабатываться с большей задержкой, например статистические отчеты и архивные журналы. Следовательно, при проектировании необходимо разделять горячие и холодные данные, использовать кэширование и выбирать подходящие модели хранения, включая БД и СУБД, соответствующие профилю конкретного сервиса.
	
	\par \redline Отдельной задачей является обеспечение отказоустойчивости. Сервис такси работает непрерывно, поэтому недоступность даже одной критической подсистемы может привести к массовым сбоям: пассажиры не смогут вызвать автомобиль, водители не получат заказы, а администрация не сможет контролировать состояние платформы. Для снижения подобных рисков применяются резервирование инстансов, балансировка нагрузки, репликация баз данных, очереди сообщений, механизмы повторной доставки событий и стратегии деградации функциональности.
	
	\par \redline Даже если на текущем этапе реализуется MVP, учет подобных факторов остается необходимым уже на стадии проектирования. Практика показывает, что игнорирование вопросов устойчивости на раннем этапе приводит к существенным трудностям при дальнейшем развитии системы. По этой причине архитектурные решения должны изначально учитывать возможность отказа отдельных компонентов и предусматривать хотя бы базовые механизмы обработки нештатных ситуаций.
	
\par \redline В высоконагруженных системах большое внимание уделяется согласованности данных. В серверной части приложения "Такси"{} одна бизнес-операция может затрагивать несколько сущностей одновременно: заказ, пользователя, водителя, платежную запись, бонусный счет и журнал событий. Если на уровне архитектуры не предусмотрены механизмы координации, то система может попасть в противоречивое состояние. Например, заказ может быть отмечен как завершенный, но платеж не будет зафиксирован, либо несколько водителей получат один и тот же заказ. Поэтому необходимо заранее определить, где требуется строгая согласованность, а где допустима временная рассогласованность данных.

\par \redline Высокая нагрузка напрямую связана с конкуренцией за ресурсы. Несколько потоков обработки могут одновременно изменять одни и те же записи, например при назначении водителя на заказ или при перерасчете стоимости поездки. Если такие ситуации не контролировать, возникают гонки данных, двойные списания средств, повторная отправка уведомлений и другие ошибки, которые негативно влияют на бизнес-процессы. Для предотвращения подобных проблем применяются оптимистические и пессимистические блокировки, идемпотентность запросов, уникальные ограничения на уровне БД и использование брокеров сообщений для упорядочивания отдельных операций.

\par \redline Важной характеристикой высоконагруженной системы является наблюдаемость. Разработчику недостаточно просто развернуть сервис и убедиться, что он принимает запросы. Необходимо собирать метрики производительности, журналы событий, трассировки распределенных запросов и показатели ошибок. В приложении такси это позволяет обнаруживать узкие места, анализировать задержки в обработке заказов, отслеживать сбои внешних интеграций и оперативно реагировать на деградацию системы.
	
	\par \redline Наблюдаемость важна не только для промышленной эксплуатации, но и для этапа разработки. Наличие журналов и базовых метрик позволяет выявлять ошибки бизнес-логики, анализировать последовательность переходов заказа между состояниями и контролировать корректность межсервисного взаимодействия. Следовательно, средства логирования и мониторинга следует рассматривать как составную часть архитектуры, а не как дополнительный элемент, вводимый после завершения основной реализации.
	
\par \redline Особое внимание должно уделяться безопасности. Серверная часть обрабатывает персональные данные пользователей, сведения о поездках, контактную информацию водителей, документы, платежные токены и историю операций. Соответственно, необходимы аутентификация и авторизация, защита каналов передачи данных, разграничение прав доступа, аудит критичных действий и механизм безопасного хранения секретов. Для внешних API также должны быть предусмотрены ограничение частоты запросов, защита от автоматизированных атак и валидация входных данных.

\par \redline Еще одна особенность высоконагруженных систем состоит в том, что не все операции должны выполняться синхронно. Некоторые действия допустимо выносить в асинхронную обработку, если это не ухудшает пользовательский сценарий. К таким операциям относятся отправка чеков, уведомлений, построение аналитических витрин, генерация отчетов, индексирование данных для поиска и часть интеграционных обменов. Использование очередей сообщений и событийного взаимодействия снижает пиковую нагрузку на центральные компоненты и делает систему более устойчивой к временным сбоям.

\par \redline Для предметной области такси критически важна также работа со временем отклика. Пользователь ожидает, что приложение покажет стоимость поездки за короткий интервал времени, а поиск водителя завершится без заметной задержки. Слишком долгая обработка воспринимается как ошибка сервиса и приводит к потере заказов. Поэтому архитектурные решения должны оцениваться не только по функциональной полноте, но и по их влиянию на скорость обработки запросов.

\par \redline Таким образом, разработка высоконагруженной серверной системы для приложения "Такси"{} требует учета масштабируемости, отказоустойчивости, согласованности данных, производительности, безопасности и наблюдаемости. Эти требования непосредственно влияют на выбор архитектурного стиля и объясняют необходимость анализа различных подходов к организации серверной части.
	
	\par
}

\subtitlespace

\subsection*{
	\gostTitleFont
	\redline
	\thechaptercntr .\thesubchaptercntr \spc
	Сравнительный анализ монолитной и микросервисной архитектур
} \addtocounter{subchaptercntr}{1}

\subtitlespace

{\gostFont
	
	\par \redline Архитектура серверного приложения определяет способ разбиения системы на логические части, характер их взаимодействия, принципы развертывания и сопровождения. Для разработки серверной части приложения "Такси"{} наиболее показательным является сравнение двух широко используемых подходов: монолитной и микросервисной архитектур [1,4].
	
	\par \redline Монолитная архитектура предполагает создание единого приложения, в котором все основные функции системы реализованы в рамках общего кода, единого процесса выполнения и, как правило, общей базы данных. В таком приложении модули могут быть логически разделены по пакетам и слоям, однако с точки зрения развертывания и масштабирования они остаются частью одного исполняемого сервиса.
	
	\par \redline Преимущество монолитного подхода состоит в относительной простоте начальной разработки. Единый проект легче собрать, протестировать и запустить, особенно на ранней стадии, когда доменная область еще уточняется. Внутренние вызовы между модулями выполняются без сетевых затрат, транзакции проще реализуются в пределах одного процесса, а требования к инфраструктуре и сопровождению обычно ниже. Для небольших систем с ограниченной функциональностью монолит может быть рациональным выбором.
	
	\par \redline Однако по мере роста приложения монолитная архитектура начинает проявлять существенные ограничения. Усложняется кодовая база, увеличивается связность между модулями, становится труднее изолировать изменения и безопасно выпускать новые версии. Масштабирование выполняется целиком для всего приложения, даже если повышенная нагрузка характерна только для одной части, например для модуля геолокации или обработки заказов. Ошибка в одном компоненте способна повлиять на доступность всей системы, а длительное время сборки и тестирования замедляет развитие продукта.
	
\par \redline Микросервисная архитектура предлагает иной подход, при котором система разбивается на набор относительно независимых сервисов, каждый из которых реализует ограниченную бизнес-область, владеет своей логикой и обычно имеет собственное хранилище данных. Такие сервисы взаимодействуют через сетевые протоколы, например HTTP, REST API или брокер сообщений, а развертывание и масштабирование могут выполняться независимо друг от друга.

\par \redline Для предметной области такси микросервисный подход выглядит естественным, поскольку бизнес-функции достаточно отчетливо разделяются. Отдельно можно выделить сервис пользователей, сервис водителей, сервис заказов, сервис тарифов, сервис платежей, сервис уведомлений, сервис геолокации и сервис аналитики. Каждый из них имеет собственный жизненный цикл, собственные требования к производительности и разные внешние интеграции. Такое разделение способствует более четкому управлению ответственностью и снижает связанность компонентов.

\par \redline К преимуществам микросервисной архитектуры относится возможность независимого масштабирования. Если в системе возрастает число запросов на расчет маршрутов и подбор водителей, можно увеличить число экземпляров только соответствующих сервисов, не затрагивая, например, административный модуль или сервис истории поездок. Кроме того, становится проще локализовать сбои, так как отказ одного сервиса не обязательно приводит к полной недоступности всей платформы. При корректной настройке механизмов деградации и повторов система может продолжить работу в частично ограниченном режиме.

\par \redline Еще одно важное преимущество микросервисов заключается в технологической и организационной гибкости. Для Java-проектов этот подход особенно удобен благодаря развитой экосистеме средств разработки и развертывания, включая JDK, Spring Boot и REST API [3]. В рамках данной работы это важно, так как проект реализуется на языке Java и должен опираться на зрелые технологические решения.
	
\par \redline Вместе с тем микросервисная архитектура влечет за собой новые сложности. Межсервисное взаимодействие происходит по сети, а значит появляются задержки, сетевые ошибки, необходимость версионирования API и контроля форматов сообщений. Транзакции между несколькими сервисами становятся сложнее, так как невозможно без дополнительных механизмов использовать привычную модель единой локальной транзакции. Возрастает роль централизованного логирования, трассировки, управления конфигурацией и балансировки нагрузки.

\par \redline Кроме того, микросервисный подход требует более тщательного проектирования границ ответственности. Если границы сервисов выбраны неудачно, система может получить избыточное число сетевых вызовов, дублирование данных и чрезмерно сложные сценарии согласования изменений. Поэтому при декомпозиции серверной части необходимо стремиться к такому разделению, при котором каждый сервис решает ограниченный набор задач и обладает четко определенным набором данных и операций.

\par \redline Существенным недостатком микросервисного подхода является повышенная сложность эксплуатации. Для стабильной работы распределенной системы требуется развитая инфраструктура развертывания, контейнеризация, автоматизация сборки и поставки, централизованный сбор метрик и продуманная политика отказоустойчивости. Если эти вопросы не решены, система может оказаться сложнее в сопровождении, чем монолит аналогичной функциональности. Следовательно, микросервисы оправданы не сами по себе, а тогда, когда их преимущества действительно покрывают инфраструктурные издержки.

\par \redline Сравнение архитектурных подходов показывает, что монолитная архитектура обеспечивает более простую начальную реализацию, тогда как микросервисная архитектура лучше соответствует требованиям масштабирования, изоляции отказов и независимого развития функциональных компонентов. Для предметной области такси данные свойства сохраняют значимость, однако на раннем этапе разработки необходимо учитывать баланс между архитектурной гибкостью и сложностью реализации.

\par \redline Для разрабатываемой серверной части приложения "Такси"{} микросервисная архитектура рассматривается как предпочтительное направление проектирования. Такой выбор обусловлен несколькими причинами. Во-первых, предметная область естественным образом разбивается на относительно автономные контексты. Во-вторых, даже минимально жизнеспособный продукт целесообразно строить с учетом возможного роста нагрузки на отдельные сервисы. В-третьих, выбранный подход создает основу для последующего расширения функциональности без полной переработки архитектуры. Наконец, использование Java и экосистемы современных фреймворков позволяет реализовать такой подход в рамках технологически обоснованного решения.

\par \redline При этом выбор микросервисной архитектуры не отменяет необходимости соблюдения инженерной дисциплины. Поэтому сервисы должны быть спроектированы с четкими границами ответственности, а межсервисные контракты - оставаться стабильными и минимально связанными.
	
	\par
}

\subtitlespace

\subsection*{
	\gostTitleFont
	\redline
	\thechaptercntr .\thesubchaptercntr \spc
	Анализ существующих решений на рынке такси-сервисов
} \addtocounter{subchaptercntr}{1}

\subtitlespace

{\gostFont
	
	\par \redline Разработка серверной части приложения "Такси"{} относится к классу задач, для которых особенно важны производительность, отказоустойчивость, масштабируемость и возможность обработки большого количества одновременных запросов. Подобные требования характерны для высоконагруженных информационных систем, где ошибки проектирования способны привести к снижению качества обслуживания пользователей и нарушению работы ключевых бизнес-процессов.
	
	\par \redline Для реализации серверной части могут использоваться различные языки программирования и платформы. Наиболее распространенными решениями являются Java, Kotlin, C\#, JavaScript и Go. Каждый из перечисленных языков обладает собственными преимуществами, однако при разработке корпоративных информационных систем необходимо учитывать зрелость экосистемы, наличие средств интеграции с базами данных, поддержку тестирования, безопасность и возможности взаимодействия с внешними сервисами.
	
	\par \redline Язык Java традиционно занимает одну из ведущих позиций в области серверной разработки благодаря стабильности, переносимости и развитой экосистеме. Важным преимуществом является наличие большого количества готовых решений для построения многослойной архитектуры, а также широкий набор библиотек и инструментов для создания масштабируемых приложений. Переносимость байт-кода между различными платформами и постоянное развитие JDK делают Java удобным инструментом для создания современных серверных систем.
	
	\par \redline Среди платформ для серверной разработки на Java особое место занимает экосистема Spring. Она позволяет реализовать четкое разделение ответственности между компонентами приложения, использовать механизмы внедрения зависимостей, организовывать REST API и упрощать работу с транзакциями, безопасностью и конфигурацией системы. Благодаря этому уменьшается объем рутинного кода и повышается сопровождаемость проекта.
	
	\par \redline В сравнении с Java другие технологии обладают определенными достоинствами. JavaScript обеспечивает высокую скорость разработки и удобен для прототипирования, однако не всегда предоставляет такой же уровень структурированности для крупных корпоративных проектов. Go демонстрирует высокую производительность и подходит для создания компактных сервисов, однако количество прикладных библиотек и готовых решений для типовых задач корпоративных информационных систем обычно уступает экосистеме Java. Платформа .NET и язык C\# также широко применяются для серверной разработки и обладают развитым набором инструментов для создания web-приложений, однако в рамках рассматриваемого проекта использование Java выглядит более целесообразным благодаря широкому распространению Spring и большому количеству накопленных практических решений.
	
	\par \redline В качестве альтернативы Java может рассматриваться Kotlin, который обладает более лаконичным синтаксисом и полной совместимостью с JVM. Тем не менее для разработки серверной части приложения "Такси"{} выбор Java является более оправданным, поскольку для данной технологии существует значительное количество документации, учебных материалов и проверенных практик построения микросервисных систем.
	
	\par \redline При выборе технологического стека необходимо учитывать не только язык программирования, но и средства разработки, позволяющие обеспечить быстрое создание приложения, удобство сопровождения и возможность последующего масштабирования. В качестве основной платформы разработки выбран Spring Boot. Данная технология ориентирована на сокращение объема конфигурации и ускорение запуска проектов. Spring Boot предоставляет набор готовых настроек для типовых сценариев, позволяет создавать автономные приложения и упрощает их развертывание в контейнерной среде.
	
	\par \redline Для серверной части приложения "Такси"{} использование Spring Boot является особенно важным, поскольку система должна обеспечивать обработку запросов пользователей, управление заказами, авторизацию, взаимодействие с внешними сервисами и обмен данными в формате JSON. Платформа предоставляет удобные средства для организации контроллеров, сервисного слоя и уровня доступа к данным, что соответствует принципам слоистой архитектуры.
	
	\par \redline Дополнительным преимуществом Spring Boot является возможность интеграции со средствами логирования, мониторинга и управления конфигурацией. Это позволяет сосредоточить основные усилия на реализации бизнес-логики приложения, а не на разработке инфраструктурных компонентов. Существенную роль играет и встроенная поддержка REST API, обеспечивающая удобный обмен данными между серверной частью и мобильными клиентами посредством HTTP-запросов.
	
	\par \redline Особое значение для приложения "Такси"{} имеет безопасность. Система должна обеспечивать защиту учетных данных пользователей, ограничение доступа к ресурсам и поддержку современных механизмов аутентификации и авторизации. Экосистема Spring предоставляет готовые инструменты для реализации данных задач, включая поддержку JWT и разграничение прав доступа.
	
	\par \redline Практическая ценность Spring Boot также проявляется в поддержке принципа разделения приложения на независимые уровни. Контроллеры отвечают за обработку HTTP-запросов, сервисный слой содержит бизнес-логику, а уровень доступа к данным обеспечивает взаимодействие с хранилищем информации. Такой подход соответствует принципам чистой архитектуры и упрощает тестирование отдельных компонентов системы.
	
	\par \redline Важным аспектом проектирования серверной части является выбор системы хранения данных. Приложение должно хранить сведения о пользователях, заказах, маршрутах, платежах и истории поездок. Для подобных задач наилучшим образом подходят реляционные системы управления базами данных, поскольку они обеспечивают поддержку транзакций, контроль целостности данных и возможность выполнения сложных запросов.
	
	\par \redline В качестве основной СУБД выбрана PostgreSQL. Данная система поддерживает транзакционность, развитые механизмы индексации и расширяемые типы данных, что делает ее удобной для хранения взаимосвязанных бизнес-сущностей. Для приложения "Такси"{} это особенно важно, поскольку информация о пользователях, водителях, заказах и платежах должна сохраняться согласованно и без нарушения связей между объектами системы.
	
	\par \redline Преимуществом PostgreSQL является также возможность выполнения сложных аналитических запросов и гибкой организации структуры данных. Это позволяет эффективно реализовывать фильтрацию поездок, поиск истории заказов, формирование отчетов и анализ пользовательской активности. Дополнительную ценность представляет поддержка расширенных механизмов индексации, которые позволяют быстро обрабатывать большие объемы информации и обеспечивать высокую скорость поиска данных.
	
	\par \redline В качестве альтернативы для хранения данных могут рассматриваться документно-ориентированные и key-value системы управления базами данных. Однако для приложения "Такси"{} подобный выбор является менее рациональным, поскольку основная информация в системе тесно связана между собой и требует строгого контроля целостности. Именно поэтому реляционная модель данных оказывается предпочтительной для реализации ядра серверной части.
	
	\par \redline Однако использование только базы данных не позволяет гарантировать минимальное время отклика системы при высокой нагрузке. В приложении такси значительная часть информации запрашивается существенно чаще, чем изменяется. К таким данным относятся настройки тарифов, справочная информация, временные состояния заказов и результаты часто повторяющихся операций. Для снижения нагрузки на основную базу данных целесообразно использовать механизм кэширования.
	
	\par \redline В качестве средства кэширования выбран Redis. Данное решение обеспечивает хранение данных в оперативной памяти и характеризуется высокой скоростью доступа. Redis поддерживает различные структуры данных и эффективно используется в сценариях, где требуется минимальная задержка при обработке запросов.
	
	\par \redline В рамках приложения "Такси"{} Redis может применяться для хранения временных состояний заказов, ограничений на повторные запросы, промежуточных результатов вычислений, пользовательских сессий и координат объектов, которые обновляются с высокой частотой. Использование Redis позволяет существенно снизить нагрузку на PostgreSQL и повысить общую производительность системы в периоды пикового спроса.
	
	\par \redline Кроме традиционного кэширования Redis может использоваться для хранения временных данных, связанных с поиском ближайшего водителя, распределением заказов между исполнителями и ограничением количества запросов от отдельных пользователей. Подобный подход способствует повышению устойчивости системы к резким скачкам нагрузки и улучшает скорость выполнения наиболее востребованных операций.
	
	\par \redline Следует отметить, что PostgreSQL и Redis решают разные задачи и не являются взаимозаменяемыми технологиями. PostgreSQL обеспечивает надежное долговременное хранение данных и сохранение их целостности, тогда как Redis выступает в роли быстрого промежуточного слоя доступа к часто используемой информации. Совместное использование данных технологий позволяет достичь оптимального баланса между надежностью хранения данных и скоростью обработки запросов.
	
	\par \redline При анализе используемых технологий также необходимо учитывать средства сборки, управления зависимостями и автоматизированного тестирования. Для Java-проектов данные задачи решаются зрелыми инструментами, которые хорошо интегрируются со Spring Boot и позволяют организовать воспроизводимую сборку приложения, централизованное управление конфигурацией и контроль качества программного кода. Это особенно важно в проектах, построенных на основе нескольких сервисов, взаимодействующих через единые программные интерфейсы.
	
	\par \redline Таким образом, для разработки серверной части приложения "Такси"{} выбран стек технологий, включающий язык программирования Java, платформу Spring Boot, систему управления базами данных PostgreSQL и средство кэширования Redis. Данный набор технологий обеспечивает высокий уровень технологической зрелости, удобство сопровождения, возможность масштабирования и соответствует требованиям, предъявляемым к современным серверным информационным системам.
	
}
	
	\subtitlespace
	
	\subsection*{
		\gostTitleFont
		\redline
		\thechaptercntr .\thesubchaptercntr \spc
		Требования к серверной части приложения
	} \addtocounter{subchaptercntr}{1}
	
	\subtitlespace
	
	{\gostFont
		
		\par \redline На основании рассмотренной предметной области и особенностей высоконагруженных систем можно сформулировать требования к серверной части приложения "Такси"{}. Эти требования определяют состав разрабатываемого решения, его архитектурные ограничения и критерии качества. При этом практическая реализация в рамках данной работы ориентирована на минимально жизнеспособный продукт, покрывающий ключевые пользовательские сценарии.
		
\par \redline Прежде всего серверная часть должна обеспечивать базовые функциональные возможности. К ним относятся регистрация и аутентификация пользователей, хранение профилей пассажиров и водителей, создание заказа, назначение водителя на поездку, сопровождение заказа по основным состояниям, расчет предварительной стоимости, обработка отмен и предоставление пользователю актуального статуса поездки.

\par \redline Указанный набор функций следует рассматривать как минимально необходимый для демонстрации работоспособности системы. Реализация именно этих сценариев позволяет подтвердить корректность выбранной архитектуры, проверить взаимодействие между ключевыми сервисами и обеспечить завершенный пользовательский процесс от создания заказа до получения результата. Функции, связанные с бонусными программами, расширенной аналитикой и сложными финансовыми операциями, могут быть отнесены к следующим этапам развития продукта.

\par \redline Для корректной работы в предметной области такси система должна поддерживать работу с геоданными. Это включает прием координат от мобильных клиентов, хранение актуального местоположения водителей, поиск ближайших исполнителей, взаимодействие с картографическими сервисами и передачу пользователям информации о маршруте и времени подачи автомобиля. Следовательно, серверная часть должна предоставлять API для оперативного обмена географическими данными и учитывать особенности их высокой изменчивости.

\par \redline Существенным функциональным требованием является поддержка ролевой модели доступа. Пассажир, водитель и администратор работают с разными сценариями, поэтому сервер обязан разграничивать права доступа и проверять полномочия на уровне каждого защищенного ресурса. Отдельно следует предусмотреть управление административными операциями, связанными с тарифами, пользователями, справочниками и мониторингом состояния системы, а также защиту соответствующих API.

\par \redline К числу нефункциональных требований относится производительность. Серверная часть должна обеспечивать приемлемое время отклика для основных пользовательских операций. Наиболее чувствительными сценариями являются аутентификация, создание заказа, поиск водителя и получение текущего статуса поездки.

\par \redline Другим важным нефункциональным требованием является масштабируемость. Так как количество пользователей и интенсивность заказов могут изменяться во времени, проектируемая система должна допускать развитие без полной переработки архитектуры. Это означает необходимость выделения независимых сервисов и использования решений, которые позволяют в дальнейшем масштабировать отдельные компоненты по мере роста нагрузки.

\par \redline Серверная часть также должна учитывать требования устойчивости к ошибкам. Даже в минимально жизнеспособном продукте необходимо предусмотреть обработку ошибок внешних интеграций, повторные попытки выполнения отдельных операций и сохранение критически важных данных. Наибольшее значение эти меры имеют для сервисов, связанных с заказами и авторизацией.

\par \redline Требования безопасности включают защищенную аутентификацию, хранение паролей в безопасном виде, использование токенов доступа, валидацию входных данных, ограничение частоты запросов и журналирование критически важных действий. Так как система работает с персональными и финансовыми данными, нарушение требований безопасности недопустимо и должно рассматриваться как один из основных рисков проекта.

\par \redline Важным требованием является сопровождаемость. Разрабатываемая система должна быть организована таким образом, чтобы внесение изменений в один функциональный блок не требовало глубокой переработки остальных компонентов. Это особенно актуально для продукта, который в дальнейшем может расширяться новыми функциями, включая бонусную систему, корпоративные тарифы, планирование поездок и интеграцию с внешними платформами [4].

\par \redline С учетом выбранного архитектурного подхода к системе можно сформулировать и инфраструктурные требования. Серверная часть должна поддерживать контейнеризированное развертывание, конфигурирование через переменные окружения и базовые средства логирования. Для межсервисного взаимодействия должны использоваться понятные контракты на основе API, а для хранения данных - решения, соответствующие профилю конкретного сервиса.
		
		\par \redline Дополнительно следует учитывать требования к удобству локального развертывания и тестирования. Для учебного проекта и MVP особенно важно, чтобы основные компоненты системы могли быть запущены в контролируемой среде без сложной ручной настройки. Это упрощает проверку работоспособности решения, ускоряет внесение изменений и делает архитектуру более прозрачной для последующего анализа в следующих разделах.
		
		\par \redline Исходя из перечисленных требований, задача данной работы состоит в разработке минимально жизнеспособного продукта серверной части приложения заказа такси, реализованного на языке Java и построенного на основе микросервисной архитектуры. Разрабатываемое решение должно обеспечивать обработку ключевых бизнес-сценариев предметной области, поддерживать взаимодействие между основными участниками сервиса и создавать основу для дальнейшего развития системы.
		
		\par \redline Для достижения поставленной цели необходимо решить следующие подзадачи:
		
		\par \redline 1. Проанализировать предметную область сервисов заказа такси и выделить основные бизнес-сущности и процессы.
		\par \redline 2. Рассмотреть особенности проектирования высоконагруженных серверных систем и определить значимые нефункциональные требования.
		\par \redline 3. Выполнить сравнительный анализ монолитной и микросервисной архитектур и обосновать выбор микросервисного подхода.
		\par \redline 4. Определить состав основных микросервисов серверной части и границы их ответственности.
		\par \redline 5. Спроектировать взаимодействие сервисов, механизмы хранения данных и способы интеграции с внешними компонентами.
		\par \redline 6. Реализовать минимально жизнеспособный продукт серверной части приложения на Java с учетом требований производительности, безопасности и дальнейшего расширения.
		
		\par \redline Таким образом, в первом разделе были определены особенности предметной области, выявлены технические ограничения и сформулированы требования, которым должен соответствовать продукт серверной части приложения "Такси"{}. Однако перед переходом к проектированию необходимо также рассмотреть существующие решения на рынке и обосновать целесообразность собственной разработки.
		
		\par
	}
	
	\subtitlespace
	
	\subsection*{
		\gostTitleFont
		\redline
		\thechaptercntr .\thesubchaptercntr \spc
		Постановка задачи
	} \addtocounter{subchaptercntr}{1}
	
	\subtitlespace
	
	{\gostFont
		
		\par \redline На сегодняшний день рынок такси-сервисов представлен как крупными международными платформами, так и локальными решениями. Среди международных игроков наиболее известны Uber, Bolt и сервисы компании Lyft. В странах СНГ сильные позиции занимает Яндекс.Такси. Каждая из этих платформ имеет свои сильные и слабые стороны, анализ которых необходим для понимания места разрабатываемого продукта на рынке.
		
		\par \redline Uber является крупнейшим в мире агрегатором такси, работающим более чем в 70 странах. Платформа предоставляет широкий спектр услуг: от классических поездок до доставки еды. Uber разработал собственный алгоритм динамического ценообразования, доказавший свою эффективность в условиях пиковых нагрузок. Однако основным недостатком Uber с точки зрения локальных рынков является зависимость от глобальной инфраструктуры: платформа использует собственные платёжные шлюзы, картографические сервисы и системы идентификации, которые в случае введения санкций или ограничений могут перестать работать на определённой территории.
		
		\par \redline Яндекс.Такси занимает лидирующие позиции в странах СНГ и Восточной Европы. Платформа глубоко интегрирована с экосистемой Яндекса: картами, навигацией, платёжной системой. Это даёт преимущество в точности маршрутизации и удобстве оплаты для пользователей из этих стран. Однако привязка к экосистеме Яндекса также является ограничением: API Яндекс.Такси закрыт для внешних разработчиков, что делает невозможным создание собственного клиентского приложения на базе этой платформы.
		
		\par \redline Bolt - европейский агрегатор такси, активно расширяющий присутствие в Восточной Европе. Платформа делает акцент на более низких комиссиях для водителей, что позволяет удерживать конкурентоспособные цены. Как и Uber, Bolt использует глобальную инфраструктуру, что создаёт те же риски в случае введения ограничений.
		
		\par \redline Общей проблемой перечисленных платформ является их зависимость от внешнеполитической ситуации. В условиях санкционного давления, введённого в последние годы, многие западные сервисы ограничили доступ к своим API, прекратили приём платежей через отдельные платёжные системы или полностью приостановили работу в некоторых странах. Это создало потребность в локальных решениях, не зависящих от внешних ограничений. Собственная платформа позволяет независимо выбирать платёжные шлюзы, картографические сервисы, центры обработки данных и прочие компоненты инфраструктуры.
		
		\par \redline Кроме политических факторов, существует и технологическая мотивация для разработки собственного решения. Коммерческие агрегаторы не предоставляют доступа к своему исходному коду и API для создания независимых клиентов. Это делает невозможным внедрение специфических функций, которые могут потребоваться локальному заказчику: особые тарифные планы, интеграция с корпоративными системами учёта, нестандартные сценарии назначения водителей или интеграция с местными картографическими сервисами. Разработка собственной платформы снимает все эти ограничения.
		
		\par \redline Таким образом, несмотря на наличие крупных игроков на рынке такси-сервисов, разработка собственного решения остаётся актуальной задачей. Собственная платформа не подвержена рискам, связанным с санкционными ограничениями, и может быть полностью адаптирована под требования локального рынка и конкретного заказчика. Это обосновывает целесообразность выполнения данного дипломного проекта.
		
		\par
	}
	
	\setcounter{subchaptercntr}{1}
	\setcounter{imagecntr}{1}
